\documentclass[11pt,a4paper]{article}
\usepackage[T1]{fontenc}
\usepackage{lmodern,amsmath,amssymb,amsthm,mathtools,booktabs,microtype}
\usepackage[margin=27mm]{geometry}
\usepackage[colorlinks=true,linkcolor=blue,urlcolor=blue,citecolor=blue]{hyperref}
\newtheorem{theorem}{Theorem}[section]
\newtheorem{lemma}[theorem]{Lemma}
\newtheorem{corollary}[theorem]{Corollary}
\newtheorem{proposition}[theorem]{Proposition}
\newcommand{\Q}{\mathbb Q}\newcommand{\Z}{\mathbb Z}
\DeclareMathOperator{\Res}{Res}
\DeclareMathOperator{\Spec}{Spec}
\setlength{\emergencystretch}{2em}
\title{Weighted bounds for eliminating finite algebraic choices}
\author{}\date{9 November 2026}
\begin{document}\maketitle
\begin{abstract}
We give a one-coordinate polynomial presentation of the rational points
of a guarded finite free algebra with reduced active fibers.
A direct scaling of its characteristic polynomials excludes singular
choices of separating element and all drops in degree.
Weighted principal minors yield explicit degree and coefficient bounds.
Applied to the specified elliptic constructions, the method gives a
family in twenty-seven total variables of degree less than $2^{151}$,
and a fully specified polynomial in twenty-nine total variables of
degree less than $2^{157}$ and coefficient height at most
$2^{2^{168}}$. Their assertions about defining the integers retain
exactly the arithmetic hypotheses of their sources; the algebraic
elimination theorem is unconditional.
\end{abstract}

\section{The guarded finite algebra}

Let $R=\Z[x_1,\ldots,x_b]$, and let
\[
 A=R[z_1,\ldots,z_m]/(f_1,\ldots,f_m)
\]
be a triangular monic presentation. Thus $f_i$ involves only
$z_1,\ldots,z_i$ and is monic of degree $e_i$ in $z_i$.
The standard monomials
\[
 z_1^{a_1}\cdots z_m^{a_m},\qquad 0\le a_i<e_i,
\]
are an $R$-basis of $A$. Write $n=\prod_i e_i$ for its rank
and $M_a$ for multiplication by $a\in A$ in this basis.
Fix a guard $G\in A$. At every rational base point we assume that
each geometric point of the fiber with $G\ne0$ has local length one.
No assumption is made about points where $G=0$; those local algebras
may be nonreduced. The predicate to be presented is
\begin{equation}\label{eq:predicate}
 \text{there is an $R$-algebra specialization $A\longrightarrow\Q$
 over the given base point with $G\ne0$.}
\end{equation}
There are no additional closed residual equations in this statement.
The two applications below have precisely this form: every equation
has already been included among the monic relations, and the remaining
conditions are nonzero guards.

For $0\le\ell\le m n(n+1)/2$, put
\[
 \lambda_\ell=1+\sum_{i=1}^m\ell^i z_i,
 \qquad H_\ell(V)=\det(I-VM_{G\lambda_\ell})
                   =\sum_{k=0}^n a_kV^k.
\]
The coefficients $a_k\in R$ depend on $\ell$, and $a_0=1$.
For $1\le d\le n$, define
\[
 H_{\ell,d}(V)=\sum_{k=0}^d a_kV^k,
 \qquad g_{\ell,d}=\Res_V(H_{\ell,d},H'_{\ell,d}).
\]
The resultant is the Sylvester determinant with prescribed degrees
$d,d-1$, including after specialization. Define
\begin{equation}\label{eq:factor}
 Q_{\ell,d}(\mathbf x,U)=
 \begin{cases}
 H_{\ell,d}(g_{\ell,d}U)^2+\displaystyle\sum_{k=d+1}^n a_k^2,
       & d<n,\\[2mm]
 H_{\ell,n}(g_{\ell,n}U),&d=n,
 \end{cases}
\end{equation}
and set
\begin{equation}\label{eq:product}
 P(\mathbf x,U)=
 \prod_{\ell=0}^{m n(n+1)/2}\ \prod_{d=1}^n Q_{\ell,d}(\mathbf x,U).
\end{equation}
All these expressions are integer polynomials. Their indices, matrices
and algebra generators are part of a finite coefficient specification,
not additional existential coordinates.

\begin{theorem}\label{thm:elimination}
For every rational base point $\mathbf x$, there is a rational $U$
with $P(\mathbf x,U)=0$ if and only if \eqref{eq:predicate} holds.
In particular, the construction adds exactly one rational coordinate.
\end{theorem}
\begin{proof}
The universal identity
\[
 \Res(f,f')=(-1)^{d(d-1)/2}a_d\operatorname{Disc}(f)
\]
shows that $g_{\ell,d}\ne0$ forces $a_d\ne0$ and makes
$H_{\ell,d}$ squarefree of degree $d$. This also holds for
$d=1$, when the discriminant is one.

Suppose $Q_{\ell,d}(\mathbf x,U)=0$. If $g_{\ell,d}=0$ or
$U=0$, then $H_{\ell,d}(g_{\ell,d}U)=1$, a contradiction.
For $d<n$, a sum of rational squares vanishes only when all its
terms vanish. Thus $a_{d+1}=\cdots=a_n=0$ and
$H_\ell=H_{\ell,d}$. The same equality is automatic for $d=n$.
Consequently $V=g_{\ell,d}U$ is a nonzero simple rational root
of $H_\ell$. With $M=M_{G\lambda_\ell}$, the rational matrix
$I-VM$ has a one-dimensional kernel. Every multiplication matrix
$M_{z_i}$ preserves this kernel and acts there by a rational scalar
$\zeta_i$. The algebra relations imply $f_i(\boldsymbol\zeta)=0$.
The eigenvalue of $M$ there is $1/V\ne0$, so
$G(\boldsymbol\zeta)\ne0$. This proves soundness, including
fibers with inactive nilpotents and complex residue fields.

Conversely, suppose a guarded rational point exists. At each geometric
point $p$ with $G(p)\ne0$, the polynomial in $\ell$
\[
 G(p)\left(1+\sum_i\ell^iz_i(p)\right)
\]
is nonzero and has degree at most $m$. For distinct such points
$p,q$, the difference of the corresponding polynomials is nonzero:
otherwise their constant coefficients give $G(p)=G(q)\ne0$,
and their other coefficients give $z_i(p)=z_i(q)$ for every $i$.
There are at most $n$ points. Counting zero values and pairwise
collisions, at most $m n(n+1)/2$ values of $\ell$ are excluded.
Hence one of the indicated integers makes all active eigenvalues
nonzero and distinct.

At an inactive local algebra multiplication by $G\lambda_\ell$
has only the eigenvalue zero, so its contribution to
$\det(I-VM)$ is one, even if the algebra is nonreduced.
All active local algebras have length one. Therefore, for the selected
$\ell$, $H_\ell$ is squarefree of some degree $1\le d\le n$.
The given rational point supplies a nonzero rational root $V$.
Its resultant $g_{\ell,d}$ is nonzero. The rational choice
$U=V/g_{\ell,d}$ annihilates the corresponding factor
in \eqref{eq:factor} and hence the product.
\end{proof}

The proof permits degree drops and does not require the algebra
or its inactive fibers to be totally real. It does not replace a
branchwise nonzero guard by the stronger requirement that its norm
be nonzero on every geometric branch.

\section{Weighted principal minors}

Give each base coordinate weight one and each $z_i$ a nonnegative
integer weight $w_i$. Assume that every nonleading monomial in
$f_i$ has weight at most $e_iw_i$. Monic reduction then never
increases weight. Suppose $G\lambda_\ell$ has weight at most
$w$, uniformly in $\ell$.

\begin{lemma}\label{lem:minors}
For the characteristic coefficients above, $\deg a_k\le kw$.
\end{lemma}
\begin{proof}
Let $b_j$ be the weight of the $j$th standard basis monomial.
The $(i,j)$ entry of $M_{G\lambda_\ell}$ has degree at most
$w+b_j-b_i$; a negative bound means that the entry is zero.
For a principal minor indexed by a set $I$ of size $k$, every
nonzero determinant term has degree at most
\[
 \sum_{j\in I}(w+b_j-b_{\sigma(j)})=kw.
\]
The basis weights cancel. Since $a_k$, up to sign, is the sum
of all principal minors of size $k$, the result follows.
\end{proof}

The resultant $\Res(H_d,H'_d)$ is weighted homogeneous of
weight $d^2$ when $a_k$ has weight $k$. Indeed, replacing $a_k$
by $s^ka_k$ replaces $H_d(V)$ by $H_d(sV)$, and the resultant
by $s^{d^2}\Res(H_d,H'_d)$. Hence
\[
 \deg g_{\ell,d}\le d^2w,
 \qquad
 \deg H_{\ell,d}(g_{\ell,d}U)
       \le (d^3+d)w+d.
\]
The degree cases can be summed separately. Write
\[
 S=1+\frac{m n(n+1)}2,
 \qquad
 C(n)=2\sum_{d=1}^{n-1}\max(d^3+d,n)+(n^3+n).
\]

\begin{theorem}\label{thm:degree}
The polynomial \eqref{eq:product} satisfies
\begin{equation}\label{eq:degreebound}
                 \deg P\le S\bigl(wC(n)+n^2\bigr).
\end{equation}
\end{theorem}
\begin{proof}
For $d<n$, the degree of its factor is at most
\[
 2\max\bigl((d^3+d)w+d,nw\bigr)
       \le2w\max(d^3+d,n)+2d.
\]
For $d=n$ the unsquared factor has degree at most
$(n^3+n)w+n$. The sum of the additive terms is
$2\sum_{d=1}^{n-1}d+n=n^2$. There are $S$ separating choices.
\end{proof}
If $t$ is the largest integer with $t^3+t<n$, then an explicit
closed formula for the finite sum is
\begin{equation}\label{eq:Cclosed}
 C(n)=\frac{n^4+3n^2}{2}+2tn
       -\frac{t^2(t+1)^2}{2}-t(t+1).
\end{equation}

\section{Coefficient norms}

For an integer polynomial $F$, write $\lVert F\rVert_1$ for
the sum of the absolute values of its coefficients. Its coefficient
height is at most $\lVert F\rVert_1$. Choose positive real weights
for the standard basis and use the resulting weighted column norm
on multiplication matrices, applying $\lVert\cdot\rVert_1$ to
each polynomial entry. Suppose
\[
                  \lVert M_{G\lambda_\ell}\rVert\le2^\beta
\]
uniformly in $\ell$. Principal minors are invariant under the
diagonal similarity defining this norm. Expanding a minor and using
column sums therefore gives
\[
 \lVert a_k\rVert_1\le {n\choose k}2^{k\beta}
                         \le2^{k(\beta+\log_2 n)}.
\]
In this section suppose $n=2^r$, and put $b=\beta+r$.

\begin{lemma}\label{lem:resnorm}
For $1\le d\le n$,
\[
 \log_2\lVert g_{\ell,d}\rVert_1
      \le d^2b+3d(r+1),
\]
where the inequality is automatic if the resultant is zero.
\end{lemma}
\begin{proof}
Every nonzero term in the Sylvester determinant has coefficient
index weight $d^2$, so the characteristic-coefficient contribution
is at most $2^{d^2b}$. There are $d$ derivative rows, contributing
at most $d^d$, and at most $(2d-1)!$ determinant terms. Their
product is at most $(2d)^{3d}$, and $\log_2(2d)\le r+1$.
\end{proof}

It follows that
\[
 \log_2\lVert H_{\ell,d}(g_{\ell,d}U)\rVert_1
       \le(d^3+d)b+3d^2(r+1)+\lceil\log_2(d+1)\rceil.
\]
For $d<n$, bounding the sum in \eqref{eq:factor} by the number
of its summands gives
\[
 \log_2\lVert Q_{\ell,d}\rVert_1
 \le2b\max(d^3+d,n)+(6d^2+3)(r+1).
\]
For $d=n$ the unsquared factor instead satisfies
\[
 \log_2\lVert Q_{\ell,n}\rVert_1
 \le(n^3+n)b+(3n^2+1)(r+1).
\]
Submultiplicativity of the coefficient norm proves the following bound.

\begin{theorem}\label{thm:norm}
Under the preceding hypotheses,
\begin{equation}\label{eq:normbound}
 \log_2\lVert P\rVert_1
 \le S\left((\beta+r)C(n)+(r+1)(2n^3+4n-2)\right).
\end{equation}
\end{theorem}
\begin{proof}
Sum the preceding inequalities over $d$ and use
\[
 6\sum_{d=1}^{n-1}d^2+3(n-1)+3n^2+1=2n^3+4n-2.
\]
Then sum over the $S$ values of $\ell$.
\end{proof}

\section{The defining-polynomial applications}

The arithmetic constructions in \cite{Source27,Source29} provide
triangular monic algebras with no closed residual equations. The
source data specify every coefficient and every nonzero guard.
They retain respectively $26$ and $28$ base coordinates, including
the distinguished input. Nine monic quartic relations describe the
point coordinates. There are respectively three and four additional
quadratic relations for pole witnesses. The point derivatives are
nonzero at every geometric point over a rational base point; the
quadratic roots themselves are guarded nonzero. Thus the active
local-length hypothesis of Theorem~\ref{thm:elimination} holds.

The following values are obtained by integer recurrences directly
from the monic relation coefficients and the guard.
\begin{center}
\begin{tabular}{lrr}
\toprule
Quantity & Twenty-seven variables & Twenty-nine variables\\
\midrule
$m$ & $12$ & $13$\\
$n$ & $2^{21}$ & $2^{22}$\\
$w$ & $6\,358\,790$ & $8\,724\,856$\\
$\beta$ at $L=137731$ & $12\,880\,956\,901$ & $12\,887\,537\,402$\\
$t$ in \eqref{eq:Cclosed} & $127$ & $161$\\
\bottomrule
\end{tabular}
\end{center}
For completeness, the generator weights for the first algebra are
\[
 \underbrace{6,\ldots,6}_{9\text{ entries}},
 1568138,\ 1935231,\ 116318,
\]
and the second has one additional weight $1935231$.
The guard weights are $4423559$ and $6789625$, respectively;
adding the largest generator weight gives $w$.
The coefficient norm bounds use the weighted companion-column
recurrences of the specified monic algebras. In each relation,
choose a logarithmic generator bound at least
\[
 \max_{k<e_i}\left\lceil
 \frac{\log_2\lVert c_{ik}\rVert}{e_i-k}\right\rceil
       +\left\lceil\log_2\#\{k:c_{ik}\ne0\}\right\rceil,
\]
where $f_i=z_i^{e_i}+\sum_{k<e_i}c_{ik}z_i^k$ and the
previously chosen generator weights are used in the coefficient
norm. This bounds multiplication by $z_i$. Applying addition and
multiplication inequalities to the guard and separating element
gives the displayed $\beta$ uniformly over its finite selector range.

\begin{corollary}\label{cor:27}
Under the arithmetic hypotheses of \cite{Source27}, there is a
computable family in $27$ total variables such that every sufficiently
large odd fixed coefficient exponent defines $\Z$ over $\Q$ and
\[
 \deg F\le
 1622838857795564438964780120527413199713157888<2^{151}.
\]
The sufficient coefficient threshold remains ineffective.
\end{corollary}
\begin{proof}
Apply Theorem~\ref{thm:elimination} to the $26$-coordinate guarded
source. Its rational base projection is preserved exactly.
The source arithmetic theorem therefore applies with the same
hypotheses and coefficient threshold. The source weights are uniform
in the odd exponent because that exponent occurs only in fixed
coefficients. Substitution into \eqref{eq:degreebound} gives the bound.
\end{proof}

\begin{corollary}\label{cor:29}
Under the arithmetic hypotheses of \cite{Source29}, the same
construction applied to its specified source with fixed exponent
$L=137731$ gives an integer polynomial in $29$ total variables
which defines $\Z$ over $\Q$, with
\[
 \deg F\le
 154383598645603405051555717167316362269845989152<2^{157},
\]
and coefficient height at most $2^{C_*}$, where
\[
 C_*=228040944767112830511621792043639838233373391751570<2^{168}.
\]
\end{corollary}
\begin{proof}
Theorem~\ref{thm:elimination} preserves the rational base projection
of the $28$-coordinate source. Its fixed parameter is the certified
sufficient one from the effective arithmetic theorem. Substitute
the second column of data into \eqref{eq:degreebound} and
\eqref{eq:normbound}.
\end{proof}

Both degree assertions are upper bounds for completely specified
polynomials, not claims that the expanded polynomials attain those
degrees. The first corollary does not certify $137731$ as a sufficient
parameter. The construction at that exponent is an algebraically
specified example in the family; its computable coefficient bound
cannot supply an effective bound for a member known to define
$\Z$. Neither corollary removes an arithmetic premise or proves
a smaller number of variables.

\section{Reproducible certificates}

The accompanying program \texttt{build\_weighted\_specification.py}
reads a verified source specification, checks the absence of closed
residuals and indexed guards, reconstructs every generator bound,
and writes \eqref{eq:product} as a complete indexed polynomial
specification. The files for the $27$- and $29$-variable constructions
include integer degree and coefficient certificates and SHA-256
identifiers for both their input and output specifications.

The independent program \texttt{check\_weighted\_specification.py}
evaluates the defining sum $C(n)$ literally and checks the integer
bounds. It also verifies small exact finite algebras with irrational
real roots, complex residue factors, zero guards, inactive nilpotents,
selector collisions, and degree drops. In every successful branch
it recovers the rational character from the one-dimensional kernel.
These tests exercise the exceptional cases in the proof; they do
not replace that proof or verify the global arithmetic hypotheses.

\begin{thebibliography}{9}
\bibitem{Source27}
\emph{A defining polynomial in twenty-seven variables:
an existence theorem with a uniform degree bound},
companion paper; its exact guarded source is supplied in
\texttt{computation/few\_variables/trace27\_nontorsion}.
\bibitem{Source29}
\emph{An explicit defining polynomial in twenty-nine variables
with effective elliptic sum heights},
companion paper; its exact guarded source is supplied in
\texttt{computation/few\_variables/effective29\_nontorsion}.
\bibitem{FiniteAlgebra}
\emph{Eliminating finite algebraic choices with one coordinate},
companion paper on general guarded finite algebras, including
closed residual equations and arbitrary number fields.
\end{thebibliography}
\end{document}
